\documentclass[12pt, a4paper]{article}

% --- Pacotes Básicos ---
\usepackage{array}
\usepackage[utf8]{inputenc}
\usepackage[T1]{fontenc}
\usepackage[main=brazilian, english]{babel}
\usepackage{indentfirst}
\usepackage{geometry}
\geometry{left=3cm, top=3cm, right=2cm, bottom=2cm}
\usepackage{booktabs}

% --- Informações do Documento ---
\title{Universidade Federal da Fronteira Sul\\
Relatório Construção de um Analisador Léxico}
\author{Erickson Giesel Müller}
\date{\today}

\begin{document}
\maketitle
\thispagestyle{empty} % primeira página não tem número.
%\newpage


\section{Introdução}
	O analisador léxico atua como a primeira fase de um compilador. Sua função é ler cada caractere do código bruto, agrupá-los em unidades lógicas com significado, chamadas lexemas, e classificá-los em categorias, chamadas tokens. O algoritmo de reconhecimento gera uma tabela de símbolos, que armazena atributos de cada lexema lido, e uma fita de saída, que representa a ordem de tokens lidos e será usada como entrada para o sintático.

	O presente trabalho tem como objetivo implementar um analisador léxico baseado no autômato finito determinístico (AFD) gerado através de uma biblioteca de tokens e gramáticas regulares. O analisador lê os caracteres em sequência e separa-os em lexemas, após o reconhecimento dos tokens, armazena na tabela de símbolos e acrescenta o identificador do token à fita de saída.

\section{Referencial Teórico}
	A função de um compilador é traduzir o código-fonte, escrito em linguagem humana, para um programa escrito em linguagem de máquina. A análise léxica separa oa palavras através dos caracteres separadores  e avalia se os caracteres digitados formam palavras que pertencem ao alfabeto da linguagem. Caso não reconheça, informa um erro para ser tratado no token.

	Para que o computador consiga reconhecer essas gramáticas, utiliza-se a Teoria dos Autômatos. Inicialmente, as regras formam um AFND, que pode conter ambiguidades e transições vazias ($\varepsilon$). Como algoritmos de reconhecimento exigem previsibilidade, aplica-se a determinização do autômato. No AFD, para cada estado e cada símbolo lido, existe apenas uma transição de destino possível.	

	A análise léxica também é responsável por lidar com entradas inválidas. Caso o algoritmo leia um caractere que não possui transição mapeada no AFD, o autômato é direcionado para um estado de erro. O analisador reporta a falha, insere o rótulo de erro na tabela de símbolos, registra a posição do lexema em que aconteceu o erro léxico e sincroniza a leitura até o próximo delimitador, permitindo que a compilação continue sem interrupções abruptas.
\section{Implementação e Resultados}
	A construção do autômato finito segue as regras determinadas a partir do arquivo \textit{fonte.txt}. A exemplo:\\

\begin{center}
	\textbf{Arquivo \textit{fonte.txt}}\\
\begin{tabular}{|>{\raggedright\arraybackslash}p{8.5cm}|} 
\toprule
se\\
sai\\
foi\\
$S::=f<A>|a<A>|e<A>|i<A>$\\
$A::=f<A>|a<A>|e<A>|i<A>|\varepsilon$\\
\bottomrule
\end{tabular}
\\
\end{center}

	Assim, temos o alfabeto $(\Sigma): \{s, e, a, i, f, o\}$ e o seguinte AFD:
\begin{center}
	\textbf{Autômato Finito Determinístico}\\
\begin{tabular}{|>{\raggedright\arraybackslash}p{1cm}|c|c|c|c|c|c|} 
    \toprule 
    & s & e & a & i & f & o \\
    \midrule 
    $\to S $   & $\{AC\}$ & I    & I    & I    & $\{FI\}$ & $\_$ \\
    $\{AC\}$   & $\_$     & B    & D    & $\_$ & $\_$     & $\_$ \\
    $*B$       & $\_$     & $\_$ & $\_$ & $\_$ & $\_$     & $\_$ \\
    $D$        & $\_$     & $\_$ & $\_$ & E    & $\_$     & $\_$ \\
    $*E$       & $\_$     & $\_$ & $\_$ & $\_$ & $\_$     & $\_$ \\
    $*\{FI\}$  & $\_$     & I    & I    & I    & I        & G    \\
    $G$        & $\_$     & $\_$ & $\_$ & H    & $\_$     & $\_$ \\
    $*H$       & $\_$     & $\_$ & $\_$ & $\_$ & $\_$     & $\_$ \\
    $*I$       & $\_$     & I    & I    & I    & I        & $\_$ \\
    $\_$       & $\_$     & $\_$ & $\_$ & $\_$ & $\_$     & $\_$ \\
    \bottomrule 
\end{tabular}
\end{center}


	A partir da função \textit{analisadorLex()}, o algoritmo de reconhecimento faz a leitura do arquivo \textit{entrada.txt}, analisando as transições de acordo com o AFD gerado. Em sua etapa de inicialização, ela carrega o conteúdo bruto do arquivo de entrada e prepara as estruturas de dados fundamentais: uma lista para armazenar a fita de saída e outra para a tabela de símbolos. Utilizando um laço de repetição de controle manual do ponteiro de leitura, o algoritmo inicia a busca por tokens sincronizando o texto. Ele ignora caracteres delimitadores prévios ($\textbackslash n$ e espaço em branco), garantindo que uma variável de contagem seja atualizada para o rastreamento preciso da localização de cada lexema no código-fonte, atualizando também a variável da linha que será guardada na tabela de símbolos.

	Após encontrar o início de uma cadeia válida, a função simula a execução do Autômato Finito Determinístico (AFD). Caractere por caractere, ela verifica as transições possíveis a partir do estado corrente. Se a transição existir, o autômato avança para o próximo estado. Caso o caractere lido não esteja no alfabeto ou force uma transição para o estado de erro, o algoritmo registra o índice exato em que a falha ocorreu, essa posição também é enviada à tabela de símbolos. Nesse cenário defensivo, o analisador continua consumindo e ignorando os caracteres da palavra inválida até encontrar um novo delimitador, impedindo que interrupções por exceção de sistema abortem a compilação.

	O índice em que o erro foi encontrado é armazenado na variável \textit{indice\_relativo}, assim o lexema pode ser recortado em duas partes, \textit{parte\_aceita} e \textit{parte\_erro}, que comporão, separadas por uma seta, a mensagem do erro na tabela de símbolos.

	Ao finalizar a leitura isolada do lexema, o algoritmo verifica se o estado alcançado pertence ao conjunto de estados de aceitação. O identificador validado é então inserido na fita de saída, que servirá de entrada para o analisador sintático. Em paralelo, a tabela de símbolos é construída incorporando os metadados do token.

	Para a demonstração prática dos resultados, utilizou-se um arquivo \textit{entrada.txt} composto por três lexemas distintos: \textit{se}, \textit{f} e \textit{fooi}. O processo de análise iniciou-se com o reconhecimento da palavra reservada \textit{se}, validada através do estado final $B$, seguida pela variável \textit{f}, reconhecida pelo estado $\{FI\}$. Esta etapa inicial confirmou a precisão do AFD em categorizar corretamente tanto cadeias fixas quanto identificadores gerados por regras gramaticais recursivas, separando-os de forma clara na tabela de símbolos.

O teste de robustez do sistema foi realizado com o lexema \textit{fooi}, que não pertence à linguagem. O analisador detectou a falha de transição no segundo caractere 'o', ativando o modo pânico e calculando o \textit{indice\_relativo} para o recorte da mensagem de erro. Como resultado, a tabela de símbolos apresentou a anotação \textit{fo→oi}, isolando visualmente o ponto exato de divergência lógica. Ao final da execução, a fita de saída e a tabela de símbolos consolidaram o sucesso do tratamento de erros e a integridade da extração de metadados necessários para as fases posteriores da compilação.
\begin{verbatim}
Fita: B F_I _ $
------------------------- TABELA DE SIMBOLOS -------------------------
Linha   | Id  | Rotulo            | Mensagem          | Palavra
______________________________________________________________________
1       | B   | palavra reservada |                   | se
______________________________________________________________________
2       | F_I | variavel          |                   | f
______________________________________________________________________
3       | _   | erro              | fo→oi             | fooi
______________________________________________________________________

\end{verbatim}

\section{Conclusão}
	O presente trabalho cumpriu com êxito o objetivo de desenvolver um analisador léxico funcional. Através da aplicação prática de conceitos de Linguagens Formais e Teoria dos Autômatos, foi possível construir um algoritmo capaz de ler dinamicamente especificações de gramáticas regulares, modelar um Autômato Finito Não Determinístico e determinizá-lo para um Autômato Finito Determinístico para o reconhecimento de tokens.

	O grande destaque da implementação está no algoritmo de leitura e no tratamento de erros léxicos. A adoção de um estado de erro, aliada à técnica de recuperação após uma transição não identificada, garantiu que o analisador não sofresse interrupções sistêmicas diante de cadeias inválidas. Além disso, a lógica desenvolvida para rastrear o índice exato de falha dentro de um lexema que deu erro, apontado como mensagem na tabela de símbolos, confere à ferramenta um caráter diagnóstico avançado, fornecendo um feedback preciso ao compilador.

	Por fim, conclui-se que as estruturas de dados geradas ao final do processo atendem plenamente aos requisitos da primeira fase de um compilador. A tabela de símbolos rica em atributos e a fita de saída limpa e validada entregam uma base sólida, isolando a complexidade textual do código-fonte e gerando os insumos exatos e necessários para alimentar as etapas subsequentes, como a análise sintática.

\begin{thebibliography}{9}
\bibitem{apostila} SCHEFFEL, Roberto M.
\textit{Apostila Linguagens Formais e Autômatos}.
\bibitem{aho}
AHO, A. V.; LAM, M. S.; SETHI, R.; ULLMAN, J. D. \textbf{Compiladores: princípios, técnicas e ferramentas}. 2. ed. São Paulo: Pearson Education do Brasil, 2007.

\bibitem{louden}
LOUDEN, K. C. \textbf{Compiladores: princípios e práticas}. São Paulo: Thomson Learning, 2004.

\end{thebibliography}

\end{document}
